%!TeX root = main.tex
%!encoding = UTF-8 Unicode
%%===========================================================================%%
%%===                          ISEA-TeX-Vorlage                           ===%%
%%===         LaTeX-Vorlage für eine schriftliche Arbeit am ISEA          ===%%
%%===		  Eine Änderungshistorie ist im GIT-Repository unter		  ===%%
%%===		 https://git.isea.rwth-aachen.de/mro-mla/latex-vorlage		  ===%%
%%===							einzusehen.								  ===%%
%%===																	  ===%%
%%===				   Vorlagenbetreuer: Michael Laumen			  		  ===%%
%%===			Anregungen/Feedback an mla@isea.rwth-aachen.de			  ===%%
%%===========================================================================%%


%%=============================================================================
%%	
%%	Die ISEA-TeX-Vorlage wurde 2007 von Matthias Bösing erstellt und von Frühjahr 2011 bis 2016 von Jens Münnix (jmu) betreut. Aktueller Betreuer (seit Mai 2016) ist Michael Laumen.
%%		
%%	
%%	Die in dieser Vorlage definierten Formatierungen sind nicht 
%%	verbindlich. Eine Anpassung des Formats an die eigenen Bedürfnisse ist ggf.
%% 	nach Rücksprache mit dem betreuenden Assistenten möglich.
%%
%%---

%%	Dokumentation zu (La)TeX findet sich lokal unter
%%		C:\Programme\MiKTeX X.Y\doc bzw. C:\Programme\MiKTeX X.Y\doc\tex\latex
%%		(der MiKTeX-Pfad ist ggf. anzupassen)
%%	Diese Vorlage verwendet das KOMA-Skript. Dokumenation:
%%		deutsch:	C:\Programme\MiKTeX X.Y\doc\latex\koma-script\scrguide.pdf
%%		englisch:	C:\Programme\MiKTeX X.Y\doc\latex\koma-script\scrguien.pdf
%%	Internetseite zum KOMA-Script-Paket mit Diskussionsforum:	
%%  	http://www.komascript.de
%%	Auf alle weiteren Fragen eine Antwort hat http://www.google.de.
%%
%%=============================================================================
\RequirePackage[utf8]{luainputenc}

%%=============================================================================
%% === Allgemeine Einstellungen ===
%%=============================================================================
%%	Diese allgemeinen Einstellungen sind vom Autor anzupassen.

%%= Sprache des Dokuments (german, english)
%% clear the precompiled files after changing the language! (tex/*.aux, tex/Main.* [except Main.tex])
\newcommand{\varLanguage}{english}

%%= Wahl des Dokumententyps (scrbook, scrartcl)
%% Hinweise: bei der Wahl von scrartcl existieren keine "chapter" und auch kein front/main/backmatter; die Vorlage ist also nicht direkt kompilierbar;
\newcommand{\varDocType}{scrbook}

%%= Doppelseitiges Dokument (true/false)
\newcommand{\varTwoside}{true}

%%= Autorenname (für Titelseite, PDF-Info und Erklärung)
\newcommand{\varAuthorGiven}{Robert}
\newcommand{\varAuthorFamily}{Kragl}
%%= Erstprüfer
\newcommand{\varSupA}{Univ.-Prof. Dr. Ingmar Kallfass}
%%= Zweitprüfer (wird nur bei Wirt.-Ing. angezeigt)
\newcommand{\varSupB}{Univ.-Prof. Dr. tbd}

%%= Datum (für Titelseite und Erklärung)
\newcommand{\varDate}{\today}

%%= Titel der Arbeit (für PDF-Info und Formblatt)
%%=	(Angabe auf Titelseite manuell, um Zeilenumbrüche selbst zu definieren)
\newcommand{\varTitle}{Measurement Environment for the EMC Characterization of Power Semiconductor Devices}

%%= Thema der Arbeit (für Titelseite, PDF-Info und Formblatt) 
\newcommand{\varSubject}{Dissertation}

%% Ist der Autor ein Wirt.-Ing? (Wichtig für Titelseite)
\newcommand{\varWirtIng}{false}

%%= Vor- oder Abgabeverision? 
%%= (bei Vorversion ist Datum in Kopfzeile und "Vorversion" mit Datum auf Titelseite)
\newcommand{\varFinalVersion}{false}

%%Dateiname der Literaturdatei (*.bib) ohne Endung
\newcommand{\varBib}{Diss}

%%= Studienrichtung
\newcommand{\varFieldOfStudy}{Electrical Engineering}

%% Institut für Logo und anpassung Declaration (Typen ISEA oder PGS)
\newcommand{\varInstitute}{ILH}

%%= 5-10 Stichworte (für PDF-Info und Formblatt)
\newcommand{\varKeywords}{Institute of Robust Power Semiconductor Systems}

%%= Wurde die Arbeit extern angefertigt? false/true (Formblatt)
\newcommand{\varExtern}{true}

%%= Bei externen Arbeiten: Firmenname der Anfertigung
\newcommand{\varExternName}{Robert Bosch GmbH}

%%= Betreuender Assistent
\newcommand{\varAssiName}{Maxime Musterfrau}



%%=============================================================================
%% === Erweiterte Einstellungen ===
%%=============================================================================
%%	Hier können die Standardeinstellungen bei Bedarf angepasst werden.

\RequirePackage{ifthen}

\ifthenelse{\equal{\varLanguage}{german}}
{
	\newcommand{\varClassLang}{ngerman}
}
{
	\newcommand{\varClassLang}{english}
}


%%=== Layout ===
%% Auswahl von Dokumentklasse und Optionen
\documentclass[%				
	BCOR=5mm,%							%% Bindekorrektur am Rand der Seite
	DIV=15,%							%% für Satzspiegel; bestimmt Größe des Textbereichs auf der Seite
	twoside=\varTwoside,%				%% s.o.
	cleardoublepage=plain,%	%% \cleardoublepage erzeugt Seite im Seitenstil 'plain'
	parskip=half,%				%% Definition von Absatzeinzug bzw. Absatzabstand
	headsepline,%					%% horizontale Linie unter Kopfzeile		
	fontsize=12pt,%				%% Schriftgröße
	listof=totoc,%				%% Eintrag ins Inhaltsverzeichnis für Abbildungs- und Tabellenverzeichnis
	index=totoc,%					%% ... für Abkürzungs- und Symbolverzeichnis
	headings=normal,%  		%% normale Größe der Überschriften
	numbers=noenddot,%		%% Punkte hinter Kapitel-, Tabellen- und Abbildungsnummerierungen vermeiden
	\varClassLang			%% Spache als globale Option 
	]%
{\varDocType}						%% s.o.
%% weitere/andere Optionen siehe KOMA-Anleitung
 
 
%%===

%\usepackage{morewrites}	%% falls Fehler "no room for a new \write" kommt, hier TESTWEISE auskommentieren
		
%%=== Sprachanpassung ===
% \ifthenelse{\equal{\varLanguage}{german}}
% {
% 	\usepackage[ngerman]{babel}
% 	\RequirePackage[ngerman=ngerman-x-latest]{hyphsubst}	%% Orthographie nach neuer deutscher Rechtschreibung	
% 	\addto\extrasngerman{\sisetup{locale = DE}}				%% autoload DE locale in siunitx for ngerman
	
% }			
% {
% 	\usepackage[english]{babel}			%% Englisch
% }

\usepackage[useregional]{datetime2}

%%=== Eingabe-Kodierung & Schriftarten ===
\usepackage{microtype}
\usepackage[T1]{fontenc}
\usepackage{lmodern}

%%=== Kopf- und Fußzeilen ===
\usepackage[%
	automark,%												%% automatische Aktualisierung der Inhalte von Kopf- und Fußzeile
	autooneside%											%% Argumente für doppelseitige Dokumente werden bei einseitigen ignoriert
]{scrlayer-scrpage}													%% Paket zur Kontrolle über Format und Inhalt von Kopf- und Fußzeile	
\pagestyle{scrheadings}							%% Auswahl des Stils von Kopf- und Fußzeilen
\automark[section]{chapter}					%% Auswahl der automatisch zu generierenden Inhalte;
	%% auf Hauptargument (hier 'chapter' = Kapitelname) wird durch \leftmark  zugegriffen;
	%% Optionales Argument (hier 'section' = Abschnittsname) wird durch \rightmark ausgewählt
\ifthenelse{\equal{\varFinalVersion}{true}}
	{\ohead{\pagemark}}								%% Seitenzahlen im Kopf aussen
	{\ohead{(\today) \quad \pagemark}}%% zusaetzlich Datum bei Vorversion
\chead{}														%% Mitte der Kopfzeile bleibt frei in Abgabeversion
\lohead{\leftmark}									%% Kapitelname rechts oben 
\rehead{\rightmark}									%% Abschnittsname links oben (nur bei doppelseitigem Druck)
\ofoot{}\cfoot{}\ifoot{}						%% Leere Fußzeile


%%===Fußnoten ===
\usepackage[multiple]{footmisc}
\renewcommand{\thefootnote}{(\roman{footnote})} %% Fußnoten im Stil: (i),(ii) etc.
%\renewcommand{\thempfootnote}{(\Roman{mpfootnote})}%% Fußnoten in Tabellen im Stil: (I),(II) etc.
%\deffootnote[Markenbreite]{Einzug}{Absatzeinzug}{Markendefinition}
\deffootnote[0.9em]{0.9em}{0.0em}{\flushbottom\textsuperscript{\thefootnotemark}}


%%=== Inhaltsverzeichnis ===
\addtocounter{secnumdepth}{+1}			%% Gliederungsebenen noch eine Ebene tiefer als voreingestellt nummerieren
\addtocounter{tocdepth}{0}					%% Ebenentiefe für Inhaltsverzeichnis bestimmen


%%=== Zeilenabstand ===
\usepackage{setspace}								%% Erlaubt die Aenderung des Zeilenabstands z.B. mittels \onehalfspacing					 %% 1 1/2 facher Zeilenabstand	

\setlength{\parindent}{0pt}			%% Kein Einzug des ersten Satzes nach einem Absatz


%%=== Mathematik ===
\usepackage{amssymb}										%% Erweiterte mathematische Formatierungen und Symbole
\usepackage{upgreek}										%% aufrechte grichische Symbole (z.B. \uppi}
\usepackage{icomma}	  									%% Abstandsanpassung hinter Dezimalkommas
% \usepackage{trsym}											%% Transformation Sybols (e.g. Vertical Laplace Sybmol)
% \usepackage{trfsigns}										%% \e Eulersche-Zahl und \j imaginäre Einheit
%\usepackage{dsfont}										%% fuer N bei Menge der nat. Zahlen (\mathds{N})
%\usepackage{esint}											%% fuer \oiint (Umlaufintegral über Linie, Fläche)

\usepackage[right]{eurosym}						%% Euro-Symbol über \euro{}; mit Betrag über \EUR{3.142}
\usepackage{textcomp}										%% erweiterte Symbole in 

\usepackage{xfrac}											%% Bruch-Funktionen (\sfrac, \dfrac, etc.)

% \usepackage[separate-uncertainty]{siunitx}					%% siunitx für Einheiten, etc.
\usepackage[uncertainty-mode=separate]{siunitx}

\ifthenelse{\equal{\varLanguage}{german}}
{
	\sisetup{quotient-mode=fraction,fraction-function=\frac,per-mode=fraction, range-phrase=\text{~bis~}}	
}			
{
	\sisetup{quotient-mode=fraction,fraction-function=\frac,per-mode=fraction, range-phrase=\text{~to~}}
}

		

\usepackage{nicefrac}


%%=== Chemie ===
\usepackage{datatool}						%% um csv (u.a.) Dateien usw einzulesen und verarbeiten zu können
\usepackage{chemfig}						%% Paket um Strukturen zu zeichnen
\usepackage[version=3,arrows=pgf]{mhchem} 	%% Paket um Strukturen und Formeln zu zeichnen
%\usepackage[xspace]{chemmacros}			%% lädt automatisch siunitx, es gibt damit ein Problem mit unitsdef


%%=== Tabellen und Spalten ===
%\usepackage{longtable}						%% ermöglicht mehrseitige Tabellen mit gleicher Spaltenbreite
%\usepackage{supertabular}					%% ermöglicht mehrseitige Tabellen
%\usepackage{dcolumn}						%% Dezimalpunkt steht untereinander in Tabs -- "may clash with longtable package"
\usepackage{multicol}						%% Wechsel zu mehrspalt. Text innerhalb einer Seite
\usepackage{multirow}						%% multirow ist ein Paket, das den gleichnamigen Befehl definiert. Es erlaubt, in einer Spalte einer Tabelle die Felder mehrerer Zeilen zu einem zusammenzufassen. 

\usepackage{booktabs}						%% schöne Tabellen

\usepackage{tabularx}
\usepackage{tablefootnote}					%% Ermöglicht Fussnoten auch in Tabellen

%\usepackage{bchart}							%% Erstellt bar-Charts sehr einfach

\usepackage{paralist}						%% Mehr Möglichketen bei Listen

\usepackage{listings}						%% Um quellcode abzubilden

%%=== Abbildungen ===
%\usepackage{flafter}									%% sorgt dafür, dass alle Float-Grafiken immer erst nach der entsprechenden Stelle eingefügt werden.

\usepackage[section]{placeins}				%% Grafiken werden nur in der dazugehörigen section eingefügt, d.h  ein zu spätes Einfügen in der nächsten section wird verhindert \FloatBarrier an Stelle, die nicht überschritten werden soll, einfügen

% \usepackage{subfig}										%% provides support for inclusion of small, `sub', figures and tables. It simplifies the positioning, captioning and labeling of such objects within a single figure or table environment and to continue a figure 
																			%% or table across multiple pages. 

%\usepackage{sidecap} 									%% defines environments SCfigure and SCtable, analogous to figure and table, which make it easy to typeset captions sideways. Additionally, a wide  environment is defined; it allows to use the margin area, e.g., for figures  wider than \textwidth.

\usepackage[labelfont=bf]{caption}								%% provides ways to customise captions in floating environments such as figure and table. 
\usepackage[list=true]{subcaption}						%%Damit kann man Unterbeschriftunge erstellen. Die Option list=true lässt diese im jeweiligen Verzeichnis erscheinen.
\usepackage{scrhack}								%% Behebt eine Warnung zwischen float und komascript
\usepackage[abs]{overpic}							%%Damit können Bilder mit einem Raster überlagert werden um Positionierungen zu vereinfachen.

\usepackage{tikz,pgf,pgfplots}					%% Auch wenn tikz kein Zeichenprogramm ist, es macht schöne Grafiken
\pgfplotsset{compat=1.3}						%% behebt einige Fehler, die beim Rendern der Grafiken auftreten können
\usetikzlibrary{plotmarks,shapes,positioning,fit}
\pgfkeys{/pgf/number format/.cd ,use comma ,set thousands separator={.}} %Globale einstellungen für Grafiken mit tikz/pgf.
\usetikzlibrary{spy} 							%% kann bei Bedarf eingebunden werden, um Lupeneffekt bei Grafiken zu ermöglichen

\pgfplotsset{
	every axis plot post/.style={
		line join=round
	}
}


%% Parameter für Grafikgröße (tikz)
\newlength\figureheight
\newlength\figurewidth 

%% Folgende Zeile einkommentieren, um externalize zu aktivieren und alle mit tikz erstellten Grafiken als pdf auszugeben.
%% -> Zwischenspeicherung der gerenderten Bilder ermöglicht schnelleres kompilieren
% \usepgfplotslibrary{external}
% \pgfkeys{/pgf/images/include external/.code=\includegraphics{#1}}

% \usepackage{filemod}

% \ifdef{\tikzexternalize}{
% 	\tikzexternalize[prefix=imgs/]					%% Zwischenspeicher im Unterordner imgs/ des tex/-Ausgabeordners
% 	\newcommand{\includetikz}[1]{
% 	  \tikzsetnextfilename{#1}
% 	  \filemodCmp{../images/tikz/#1.tex}{imgs/#1.pdf}
% 	    {\tikzset{external/remake next}}{}
% 	  \input{../images/tikz/#1.tex}
% 	}
% }
% {
% 	\newcommand{\includetikz}[1]{
% 		\input{../images/tikz/#1.tex}
% 	}
% }

\setcounter{secnumdepth}{4} % To number paragraphs
\setcounter{tocdepth}{4}


\usepackage{import}     % For better handling of included files with relative paths
\usepackage{standalone} % For correctly embedding standalone LaTeX figures

\usepackage{tikz}
\usetikzlibrary{external}
\tikzexternalize[prefix=tikz_extern/]
\graphicspath{{./figures/}}
\makeatletter
\def\input@path{{./figures/}}




%%=== Sonstiges ===
\usepackage{graphicx}						%% \usepackage[draft]{graphicx} --> Grafiken werden nur als leere Box ausgegeben
\usepackage{color}
\usepackage{pdflscape}						%% (einzelne) Seiten koennen im Querformat ins PDF gestellt werden

\definecolor{isea_blue}{rgb}{0.031, 0.020, 0.580}	%% ISEA-Blau fuer PDF links,
	%% wird in CustomVariables.tex erneut defniniert; hier, um in PDF-links nutzbar zu sein;
	%% CustomVariables.tex wird erst in Main.tex eingebunden, um fuer den Nutzer, der eigene Eintraege 
	%% dort ergaenzen kann, erkennbar zu sein

\ifthenelse{\equal{\varLanguage}{german}}
	{%\usepackage[ngerman]{cleveref}						%%Für Kommentare und ToDo Listen
	\usepackage[ngerman]{todonotes}}
	{%\usepackage[english]{cleveref}						%%Für Kommentare und ToDo Listen
	\usepackage[english]{todonotes}}

\usepackage{pdfpages}
\usepackage{pdfcomment}

\newcommand{\mypdfcomment}[1]{\pdfcomment[author=\varAuthorGiven,markup=Highlight,color=yellow]{#1}}
\newcommand{\pdftext}[1]{\pdfannot{/Subtype/Text/Contents(#1)}}
\newcommand{\mypdfmarkupcomment}[2]{\pdfmarkupcomment[author=\varAuthorGiven{} \varAuthorFamily{},markup=Highlight,color=yellow]{#1}{#2}}
\newcommand{\mypdfmargincomment}[1]{\pdfmargincomment[author=\varAuthorGiven{} \varAuthorFamily{},markup=Highlight,color=yellow]{#1}}

%%== Anführungszeichen setzen ==
\ifthenelse{\equal{\varLanguage}{german}}
	{\usepackage[autostyle=true,german=quotes]{csquotes}}			%% Eine Ergänzung für Anführungszeichen
	{\usepackage[autostyle=true,english=british]{csquotes}}			%% Englisch
%% Doku zu csquotes: http://www.ctan.org/tex-archive/macros/latex/contrib/csquotes/csquotes.pdf

%%== Literaturverzeichnis ==
\usepackage[style=numeric-comp, minnames=1, maxnames=2, hyperref=true, backend=biber, bibencoding=utf8]{biblatex}		%% Lädt das Biblatex Paket und definert den Style


%% Styles: numeric für Zahlen in eckigen Klammern [1]

%\newcounter{biburllcpenalty} % definiere die neuen counter
%\newcounter{biburlucpenalty} % definiere die neuen counter

\apptocmd{\UrlBreaks}{\do\f\do\m}{}{} % Einstellung um URLs korrekt darzustellen (Zeilenumbruch)
\setcounter{biburllcpenalty}{9000}% Kleinbuchstaben
\setcounter{biburlucpenalty}{9000}% Großbuchstaben


%%== Literaturverzeichnis (alternativ, nutzt den Bibkey als Text ([Gru96] statt [1])
%\usepackage[style=alphabetic, minnames=1, maxnames=2, hyperref=true, backend=biber, bibencoding=utf8]{biblatex}
%\DeclareFieldFormat{labelalpha}{\thefield{entrykey}}
%\DeclareFieldFormat{extraalpha}{}
%%== Ende Literaturverzeichnis (alternativ)
	
\ifthenelse{\equal{\varLanguage}{german}}
	{\DefineBibliographyStrings{ngerman}{bibliography={Literaturverzeichnis}}
	\DefineBibliographyStrings{ngerman}{mathesis={Masterarbeit}}
	\DefineBibliographyStrings{ngerman}{phdthesis={Dissertation}}
	}% NICHT references
	{}

\makeatletter
\@ifpackagelater{biblatex}{2016/03/03} % biblatex syntax changed with version 3.3 released on 2016/03/03
{ % biblatex version >= 3.3
	% New syntax for flexible back end
	\DeclareNameFormat{default}{%
		\renewcommand*{\multinamedelim}{\addsemicolon\addspace}%
		\nameparts{#1}%
		\usebibmacro{name:family-given}
		{\namepartfamily}
		{\namepartgiveni}
		{\namepartprefixi}
		{\namepartsuffixi}%
		\usebibmacro{name:andothers}%
	}
	
	\DeclareNameFormat{labelname}{%
		\renewcommand*{\multinamedelim}{\addcomma\addspace}%
		\nameparts{#1}%
		\usebibmacro{name:family-given}
		{\namepartfamily}
		{\namepartgiveni}
		{\namepartprefixi}
		{\namepartsuffixi}%
		\usebibmacro{name:andothers}%
	}
	
	\DeclareNameFormat{sortname}{%
		\renewcommand*{\multinamedelim}{\addcomma\addspace}%
		\nameparts{#1}%
		\usebibmacro{name:family-given}
		{\namepartfamily}
		{\namepartgiveni}
		{\namepartprefixi}
		{\namepartsuffixi}%
		\usebibmacro{name:andothers}%
	}	
}
{ % biblatex version < 3.3
	%%Literaturverzeichnis mit Nachname, Vorname (abgekürzt)
	\DeclareNameFormat{default}{%
		\usebibmacro{name:last-first}{#1}{#4}{#6}{#8}%
		\usebibmacro{name:andothers}}
	\DeclareNameFormat{labelname}{%
		\usebibmacro{name:last-first}{#1}{#4}{#6}{#8}%
		\usebibmacro{name:andothers}}
	\DeclareNameFormat{sortname}{%
		\usebibmacro{name:last-first}{#1}{#4}{#6}{#8}%
		\usebibmacro{name:andothers}}
}
\makeatother

\bibliography{../literature/\varBib}		%% Pfad der Bibliotheks-Datei. Hier nichts ändern!

%%=== Index ===
%\usepackage{makeidx}
%\makeindex

%%=== PDF ===
\usepackage{ifpdf}

\newboolean{bool_hyperref}
\setboolean{bool_hyperref}{true}		%% Soll hyperref (pdf-links) geladen werden?

\ifpdf
	\ifthenelse{\boolean{bool_hyperref}}{
	    \usepackage{hyperref}   				%% Hyperlinks	(hyperref als allerletztes Paket laden!!!)
	    \hypersetup
	    {
	        pdftitle		={\varTitle},
	        pdfsubject	={\varSubject},
	        pdfauthor		={\varAuthorGiven{} \varAuthorFamily{}},
	        pdfkeywords	={\varKeywords},
	        colorlinks	=true,
	        linkcolor		=black, 
	        citecolor		=black, 
	        urlcolor		=black,
	        pdfstartview=Fit, %FitH, FitB	
	        %plainpages=false
	    }
		\ifthenelse{\equal{\varTwoside}{true}}
		{
			\hypersetup{pdfduplex=DuplexFlipLongEdge}
		}{
			\hypersetup{pdfduplex=Simplex}
		}
	}{ }
    \graphicspath{{../images/}{../images/}}
    \ifthenelse{\equal{\varFinalVersion}{true}}
    	{\pdfcompresslevel=9}    				%% höchste und langsamste Kompression
    	{\pdfcompresslevel=0}    				%% geringste und schnellste Kompression
    %\pdfhorigin	= 1cm
    %\pdfvorigin	=	1cm
    %\hoffset			=	0.7mm
    %\voffset			=	5mm
\else
    \graphicspath{{../images/}}
		%\hoffset			= 0.7mm
		%\voffset			= -0.2mm
\fi

%%=== Verzeichnisse (Symbol, Abkürzungen usw.) ===

\usepackage[toc=true,nonumberlist,acronym,acronymlists={symbolslist},nomain]{glossaries}
%% ftp://ftp.rrzn.uni-hannover.de/pub/mirror/tex-archive/macros/latex/contrib/glossaries/glossaries-user.pdf
%% ftp://ftp.dante.de/tex-archive/macros/latex/contrib/glossaries/glossariesbegin.pdf

\renewcommand*{\glspostdescription}{}		%% Schaltet den Punkt am Ende der Erklärung ab.
\newglossary[sya]{symbolslist}{syb}{syc}{Symbolverzeichnis}
%\newglossary[fra]{fremdwort}{frb}{frc}{Fremdwörterverzeichnis}

%Glossar Befehle einschalten
\makeglossaries